\documentclass[11pt]{article}

\usepackage[margin=1in]{geometry}
\usepackage{amsmath, amssymb}
\usepackage{array}
\usepackage{enumitem}

\title{Estimator Formula Audit for the Empirical Proof of Concept}
\author{}
\date{}

\begin{document}

\maketitle

\section{Purpose}

This document writes the exact estimator formulas for the empirical proof-of-concept versions of Naive Method I and Naive Method II.

The formulas below intentionally do not introduce intermediate symbols such as logit differences or log-count-ratio variables. Every displayed estimator input is either a count of an observed token or a base next-token logit.

\section{Token Pair}

The proof-of-concept experiment uses one token pair:

\begin{itemize}
    \item Token subscript \(0\) means Qwen2.5 token ID \(15\), whose decoded string is \texttt{0}.
    \item Token subscript \(1\) means Qwen2.5 token ID \(16\), whose decoded string is \texttt{1}.
\end{itemize}

Thus, for example, \(n_{2,0}\) is a count for token ID \(15\), while \(z_{2,1}\) is a logit for token ID \(16\).

\section{Prompt Inventory}

The six discovered prompts are:

\subsection*{Prompt 1}
\begin{verbatim}
Output a random binary sequence and nothing else. I should not be seeing words in your response.
\end{verbatim}

\subsection*{Prompt 2}
\begin{verbatim}
Generate a RANDOM number from {zero, one}. Your response should be a single numerical digit. Nothing else. Try to be as unpredictable as possible.
\end{verbatim}

\subsection*{Prompt 3}
\begin{verbatim}
Generate a highly random sequence of zero's and one's. Your response should be a string of numerical digits and nothing else. Do NOT output non-numerical text.
\end{verbatim}

\subsection*{Prompt 4}
\begin{verbatim}
Simulate a completely random coin flip. Output the number zero if it's heads. Output the number one if it's tails. Your output should be a single numerical digit. DON'T GIVE THE SAME RESPONSE EVERY TIME. TRY TO MAKE IT A FIFTY-FIFTY CHANCE.
\end{verbatim}

\subsection*{Prompt 5}
\begin{verbatim}
Generate a completely random comma-separated list of zero's and ones. Each entry in the list must be a numerical digit. Do not output anything else except this list.
\end{verbatim}

\subsection*{Prompt 6}
\begin{verbatim}
OUTPUT THE NUMBER ZERO OR THE NUMBER ONE. MAKE IT A FIFTY-FIFTY CHANCE. YOUR RESPONSE SHOULD BE A NUMERICAL DIGIT AND NOTHING ELSE.
\end{verbatim}

\section{Canary Context Choices}

The paper's general naive methods use four context slots \(C_1,C_2,C_3,C_4\). In this empirical version, we impose
\[
C_3 = C_2.
\]
Therefore there are only three distinct context inputs: \(C_1\), \(C_2\), and \(C_4\). The formulas in this document have already substituted \(C_3=C_2\), so no separate \(C_3\) count or logit appears.

\begin{center}
\begin{tabular}{|>{\raggedright\arraybackslash}p{0.18\linewidth}|>{\raggedright\arraybackslash}p{0.22\linewidth}|>{\raggedright\arraybackslash}p{0.22\linewidth}|>{\raggedright\arraybackslash}p{0.22\linewidth}|}
\hline
Reference model & \(C_1\) & \(C_2=C_3\) & \(C_4\) \\
\hline
Qwen2.5-7B-Instruct & Prompt 1 & Prompt 2 & Prompt 3 \\
\hline
Qwen2.5-14B-Instruct & Prompt 1 & Prompt 4 & Prompt 5 \\
\hline
Qwen2.5-32B-Instruct & Prompt 2 & Prompt 6 & Prompt 4 \\
\hline
\end{tabular}
\end{center}

\section{Inputs}

Fix one row of the canary context table above. All formulas below are for that fixed canary row.

\begin{itemize}[leftmargin=*]
    \item \(j \in \{1,2,4\}\) is the context-slot subscript. Slot \(3\) from the paper has been identified with slot \(2\), so the empirical version uses \(1,2,4\) as the three distinct context slots.
    \item \(n_{j,0}\) is the number of times the candidate black-box LLM B outputs token ID \(15\), decoded as \texttt{0}, as the first next token when sampled on context \(C_j\).
    \item \(n_{j,1}\) is the number of times the candidate black-box LLM B outputs token ID \(16\), decoded as \texttt{1}, as the first next token when sampled on context \(C_j\).
    \item \(z_{j,0}\) is the owned/reference LLM A's base next-token logit for token ID \(15\), decoded as \texttt{0}, after context \(C_j\).
    \item \(z_{j,1}\) is the owned/reference LLM A's base next-token logit for token ID \(16\), decoded as \texttt{1}, after context \(C_j\).
\end{itemize}

Since \(C_3=C_2\), the slot-3 logit inputs from the paper are already replaced by \(z_{2,0}\) and \(z_{2,1}\), and the slot-3 count inputs from the paper are already replaced by \(n_{2,0}\) and \(n_{2,1}\).

The formulas require the count inputs inside logarithms to be positive. If any of \(n_{1,0},n_{1,1},n_{2,0},n_{2,1},n_{4,0},n_{4,1}\) is zero, the finite-sample estimator value is not yet defined for that sample state.

\section{Runtime \(g\)-Inputs}

For implementation, we first compute the context-level quantities:

\[
g_1:=\ln(n_{1,0})-\ln(n_{1,1})
\]

\[
g_2:=\ln(n_{2,0})-\ln(n_{2,1})
\]

\[
g_4:=\ln(n_{4,0})-\ln(n_{4,1})
\]

By default, this means \(a=\texttt{0}\), \(b=\texttt{1}\), token \(a\) has token ID \(15\), and token \(b\) has token ID \(16\). Thus the runtime convention is:

\[
g_j=\ln(n_{j,a})-\ln(n_{j,b})
\]

with \(a=0\) and \(b=1\) by default.

The fully expanded formulas below are the audit formulas. The later runtime formulas are the same formulas after replacing each log-count difference by its corresponding \(g_j\).

\section{Our Naive Method I}

The empirical Naive Method I estimator is:

\[
\boxed{
\begin{aligned}
&
\frac{
    \left(z_{1,0}-z_{1,1}\right)-\left(z_{2,0}-z_{2,1}\right)
}{
    \left(\ln(n_{1,0})-\ln(n_{1,1})\right)
    -
    \left(\ln(n_{2,0})-\ln(n_{2,1})\right)
}
\\[1.25em]
&\quad -
\frac{
    \left(z_{2,0}-z_{2,1}\right)-\left(z_{4,0}-z_{4,1}\right)
}{
    \left(\ln(n_{2,0})-\ln(n_{2,1})\right)
    -
    \left(\ln(n_{4,0})-\ln(n_{4,1})\right)
}
\end{aligned}
}
\]

This is the paper's Naive Method I with \(k=k'=1\), the token pair \(a_1,b_1\) equal to the token pair \(a'_1,b'_1\), and \(C_3=C_2\) substituted directly into the formula.

If LLM A equals LLM B and the regularity assumptions from the paper hold, this estimator is intended to converge to \(0\).

\subsection*{Runtime \(g\)-Version}

In implementation, Naive Method I is computed from \(g_1,g_2,g_4\) as:

\[
\boxed{
\frac{
    \left(z_{1,0}-z_{1,1}\right)-\left(z_{2,0}-z_{2,1}\right)
}{
    g_1-g_2
}
-
\frac{
    \left(z_{2,0}-z_{2,1}\right)-\left(z_{4,0}-z_{4,1}\right)
}{
    g_2-g_4
}
}
\]

\section{Our Naive Method II}

The empirical Naive Method II estimator is:

\[
\boxed{
\begin{aligned}
&
\frac{
    \left(\ln(n_{2,0})-\ln(n_{2,1})\right)
    \left(z_{1,0}-z_{1,1}\right)
    -
    \left(\ln(n_{1,0})-\ln(n_{1,1})\right)
    \left(z_{2,0}-z_{2,1}\right)
}{
    \left(\ln(n_{1,0})-\ln(n_{1,1})\right)
    -
    \left(\ln(n_{2,0})-\ln(n_{2,1})\right)
}
\\[1.25em]
&\quad -
\frac{
    \left(\ln(n_{4,0})-\ln(n_{4,1})\right)
    \left(z_{2,0}-z_{2,1}\right)
    -
    \left(\ln(n_{2,0})-\ln(n_{2,1})\right)
    \left(z_{4,0}-z_{4,1}\right)
}{
    \left(\ln(n_{2,0})-\ln(n_{2,1})\right)
    -
    \left(\ln(n_{4,0})-\ln(n_{4,1})\right)
}
\end{aligned}
}
\]

This is the paper's Naive Method II with \(k=k'=1\), the same token pair used in both estimator components, and \(C_3=C_2\) substituted directly into the formula.

If LLM A equals LLM B and the regularity assumptions from the paper hold, this estimator is intended to converge to \(0\).

\subsection*{Runtime \(g\)-Version}

In implementation, Naive Method II is computed from \(g_1,g_2,g_4\) as:

\[
\boxed{
\frac{
    g_2\left(z_{1,0}-z_{1,1}\right)
    -
    g_1\left(z_{2,0}-z_{2,1}\right)
}{
    g_1-g_2
}
-
\frac{
    g_4\left(z_{2,0}-z_{2,1}\right)
    -
    g_2\left(z_{4,0}-z_{4,1}\right)
}{
    g_2-g_4
}
}
\]

\section{Implementation Checklist}

For a fixed reference model row, the implementation needs exactly the following numerical inputs:

\begin{itemize}[leftmargin=*]
    \item Context-level values from candidate LLM B samples: \(g_1,g_2,g_4\).
    \item Reference logits from owned LLM A: \(z_{1,0}, z_{1,1}, z_{2,0}, z_{2,1}, z_{4,0}, z_{4,1}\).
\end{itemize}

The raw counts \(n_{1,0}, n_{1,1}, n_{2,0}, n_{2,1}, n_{4,0}, n_{4,1}\) are used upstream to produce \(g_1,g_2,g_4\). The estimator functions themselves consume \(g\)-values rather than raw count dictionaries.

No probability inputs, logit-bias inputs, temperature inputs, or shorthand logit-difference variables are required by the runtime estimator functions.

\end{document}
